% GENERATED FILE. Edit canonical lessons, then run npm run generate:books.
% canonical-source-hash: fnv1a64:e13524aa362bf18f
% canonical-lessons: ES-C551-acogida, ES-C551-instalarse, ES-C551-establecerse, ES-C551-emigrar, ES-C551-adaptarse, ES-R551-repaso-544-551

\chapter{Una vida nueva}
\label{ch:es-una-vida-nueva}
\hlchaptermodality{551}

\begin{chapteropening}
\textbf{By the end of this chapter:} \emph{I can say that someone found a warm welcome, settled in, established themselves, emigrated and adapted.}

Complete the last lesson of chapter 551: I can say that someone found a warm welcome, settled in, established themselves, emigrated and adapted.
\end{chapteropening}

\section[la acogida]{la acogida — the welcome, the reception}

\label{lesson:ES-C551-acogida}



\begin{quote}
Before the new one: say the Spanish for the widower, then the Spanish for the widow.
\end{quote}

\subsection*{You'll want to know: la acogida}
\textbf{la acogida} — \textquotedblleft{}the welcome, the reception\textquotedblright{}.

\begin{quote}
La acogida en el pueblo fue muy buena.
\end{quote}

\begin{quote}
\emph{The welcome in the village was very good.}
\end{quote}

\emph{Un centro de acogida} is a shelter for refugees or homeless people.

\begin{cousinweb}
\emph{acogida} is built on \emph{acoger} 'to take in, to welcome' (EtymDB), which is inherited from spoken Latin \emph{*accolligere}, built on \emph{colligere} 'to gather together'.
\end{cousinweb}

\begin{grammarlens}[title={What you've built}]
One more word for telling, explaining and describing at B1.
\end{grammarlens}

\subsection*{Your turn}
\begin{itemize}
\raggedright
  \item \emph{Say it:} \emph{la acogida}
  \item \emph{Say it:} \emph{la acogida}, once more
  \item \emph{Say it:} a sentence of your own with \emph{la acogida}
  \item \emph{Recall:} say the Spanish for the widower, then the Spanish for the widow, then say \emph{la acogida} again
\end{itemize}

\begin{tcolorbox}[breakable,colback=teal!4,colframe=teal!35!black,title={Before you move on}]
What does la acogida mean? (\textquotedblleft{}The welcome, the reception\textquotedblright{}.) Say it once more.
\end{tcolorbox}
\section[instalarse]{instalarse — to settle in, to move in}

\label{lesson:ES-C551-instalarse}



\begin{quote}
Before the new one: say the Spanish for the widow, then the Spanish for the welcome.
\end{quote}

\subsection*{You'll want to know: instalarse}
\textbf{instalarse} — \textquotedblleft{}to settle in, to move in\textquotedblright{}.

\begin{quote}
Nos instalamos en un piso pequeño.
\end{quote}

\begin{quote}
\emph{We settled into a small flat.}
\end{quote}

Without \emph{se}: \emph{instalar} is to install.

\begin{cousinweb}
\emph{instalarse} was borrowed from French \emph{installer} (EtymDB), from Medieval Latin \emph{installāre} 'to put into a seat', built on \emph{stallum} 'a stall, a seat', a Germanic word.
\end{cousinweb}

\begin{grammarlens}[title={What you've built}]
One more word for telling, explaining and describing at B1.
\end{grammarlens}

\subsection*{Your turn}
\begin{itemize}
\raggedright
  \item \emph{Say it:} \emph{instalarse}
  \item \emph{Say it:} \emph{instalarse}, once more
  \item \emph{Say it:} a sentence of your own with \emph{instalarse}
  \item \emph{Recall:} say the Spanish for the widow, then the Spanish for the welcome, then say \emph{instalarse} again
\end{itemize}

\begin{tcolorbox}[breakable,colback=teal!4,colframe=teal!35!black,title={Before you move on}]
What does instalarse mean? (\textquotedblleft{}To settle in, to move in\textquotedblright{}.) Say it once more.
\end{tcolorbox}
\section[establecerse]{establecerse — to settle, to establish oneself}

\label{lesson:ES-C551-establecerse}



\begin{quote}
Before the new one: say the Spanish for the welcome, then the Spanish for to settle in.
\end{quote}

\subsection*{You'll want to know: establecerse}
\textbf{establecerse} — \textquotedblleft{}to settle, to establish oneself\textquotedblright{}.

\begin{quote}
La familia se estableció en Argentina.
\end{quote}

\begin{quote}
\emph{The family settled in Argentina.}
\end{quote}

It takes a \emph{z}: \emph{me establezco}.

\begin{cousinweb}
\emph{establecerse} comes from spoken Latin \emph{*stabilīscere} (EtymDB), built on \emph{stabilīre} 'to make firm' (Lewis \& Short), from \emph{stabilis} 'steady', built on \emph{stāre} 'to stand'.
\end{cousinweb}

\begin{grammarlens}[title={What you've built}]
One more word for telling, explaining and describing at B1.
\end{grammarlens}

\subsection*{Your turn}
\begin{itemize}
\raggedright
  \item \emph{Say it:} \emph{establecerse}
  \item \emph{Say it:} \emph{establecerse}, once more
  \item \emph{Say it:} a sentence of your own with \emph{establecerse}
  \item \emph{Recall:} say the Spanish for the welcome, then the Spanish for to settle in, then say \emph{establecerse} again
\end{itemize}

\begin{tcolorbox}[breakable,colback=teal!4,colframe=teal!35!black,title={Before you move on}]
What does establecerse mean? (\textquotedblleft{}To settle, to establish oneself\textquotedblright{}.) Say it once more.
\end{tcolorbox}
\section[emigrar]{emigrar — to emigrate}

\label{lesson:ES-C551-emigrar}



\begin{quote}
Before the new one: say the Spanish for to settle in, then the Spanish for to settle.
\end{quote}

\subsection*{You'll want to know: emigrar}
\textbf{emigrar} — \textquotedblleft{}to emigrate\textquotedblright{}.

\begin{quote}
Mis abuelos emigraron a Alemania.
\end{quote}

\begin{quote}
\emph{My grandparents emigrated to Germany.}
\end{quote}

\emph{Emigrar} is leaving your country; \emph{inmigrar} is coming into a new one.

\begin{cousinweb}
\emph{emigrar} is a learned borrowing from Latin \emph{ēmigrāre} 'to move out, to depart' (EtymDB), built on \emph{ē} 'out' and \emph{migrāre} 'to move'.
\end{cousinweb}

\begin{grammarlens}[title={What you've built}]
One more word for telling, explaining and describing at B1.
\end{grammarlens}

\subsection*{Your turn}
\begin{itemize}
\raggedright
  \item \emph{Say it:} \emph{emigrar}
  \item \emph{Say it:} \emph{emigrar}, once more
  \item \emph{Say it:} a sentence of your own with \emph{emigrar}
  \item \emph{Recall:} say the Spanish for to settle in, then the Spanish for to settle, then say \emph{emigrar} again
\end{itemize}

\begin{tcolorbox}[breakable,colback=teal!4,colframe=teal!35!black,title={Before you move on}]
What does emigrar mean? (\textquotedblleft{}To emigrate\textquotedblright{}.) Say it once more.
\end{tcolorbox}
\section[adaptarse]{adaptarse — to adapt, to get used to}

\label{lesson:ES-C551-adaptarse}



\begin{quote}
Before the new one: say the Spanish for to settle, then the Spanish for to emigrate.
\end{quote}

\subsection*{You'll want to know: adaptarse}
\textbf{adaptarse} — \textquotedblleft{}to adapt, to get used to\textquotedblright{}.

\begin{quote}
Me adapté rápido a la nueva ciudad.
\end{quote}

\begin{quote}
\emph{I quickly got used to the new city.}
\end{quote}

\emph{Adaptarse a} — to adapt to.

\begin{cousinweb}
\emph{adaptarse} is a learned borrowing from Latin \emph{adaptāre} 'to fit to' (Lewis \& Short), built on \emph{ad} and \emph{aptus} 'fitted, suited'.
\end{cousinweb}

\begin{grammarlens}[title={What you've built}]
One more word for telling, explaining and describing at B1.
\end{grammarlens}

\subsection*{Your turn}
\begin{itemize}
\raggedright
  \item \emph{Say it:} \emph{adaptarse}
  \item \emph{Say it:} \emph{adaptarse}, once more
  \item \emph{Say it:} a sentence of your own with \emph{adaptarse}
  \item \emph{Recall:} say the Spanish for to settle, then the Spanish for to emigrate, then say \emph{adaptarse} again
\end{itemize}

\begin{tcolorbox}[breakable,colback=teal!4,colframe=teal!35!black,title={Before you move on}]
What does adaptarse mean? (\textquotedblleft{}To adapt, to get used to\textquotedblright{}.) Say it once more.
\end{tcolorbox}
\section[(dialogue)]{Repaso de los capítulos 544–551}

\label{lesson:ES-R551-repaso-544-551}



\begin{quote}
Say the Spanish for the birth and the delivery.
\end{quote}

\subsection*{Your turn}
Say these aloud:

\begin{itemize}
\raggedright
  \item the birth, the delivery, to give birth, the pregnancy, the newborn baby
  \item to bring up, childhood, youth, adolescence, the generation
  \item the christening, the godfather, the godmother, the grandson, the nappy
  \item the boyfriend, the girlfriend, to get engaged, the engagement, the ring
  \item the marriage, the ceremony, the banquet, the toast, the honeymoon
  \item married, single, to marry, the husband, to toast
  \item the divorce, to separate, divorced, the widower, the widow
  \item the welcome, to settle in, to settle, to emigrate, to adapt
\end{itemize}

\begin{tcolorbox}[breakable,colback=teal!4,colframe=teal!35!black,title={Before you move on}]
The birth? (\textbf{el nacimiento}.) The delivery? (\textbf{el parto}.) To give birth? (\textbf{dar a luz}.) The pregnancy? (\textbf{el embarazo}.) The newborn baby? (\textbf{el recién nacido}.) To bring up? (\textbf{criar}.) Childhood? (\textbf{la niñez}.) Youth? (\textbf{la juventud}.) Adolescence? (\textbf{la adolescencia}.) The generation? (\textbf{la generación}.) The christening? (\textbf{el bautizo}.) The godfather? (\textbf{el padrino}.) The godmother? (\textbf{la madrina}.) The grandson? (\textbf{el nieto}.) The nappy? (\textbf{el pañal}.) The boyfriend? (\textbf{el novio}.) The girlfriend? (\textbf{la novia}.) To get engaged? (\textbf{comprometerse}.) The engagement? (\textbf{el compromiso}.) The ring? (\textbf{el anillo}.) The marriage? (\textbf{el matrimonio}.) The ceremony? (\textbf{la ceremonia}.) The banquet? (\textbf{el banquete}.) The toast? (\textbf{el brindis}.) The honeymoon? (\textbf{la luna de miel}.) Married? (\textbf{casado}.) Single? (\textbf{soltero}.) To marry? (\textbf{casar}.) The husband? (\textbf{el marido}.) To toast? (\textbf{brindar}.) The divorce? (\textbf{el divorcio}.) To separate? (\textbf{separarse}.) Divorced? (\textbf{divorciado}.) The widower? (\textbf{el viudo}.) The widow? (\textbf{la viuda}.) The welcome? (\textbf{la acogida}.) To settle in? (\textbf{instalarse}.) To settle? (\textbf{establecerse}.) To emigrate? (\textbf{emigrar}.) To adapt? (\textbf{adaptarse}.)
\end{tcolorbox}
